\chapter{Regular Languages and Lexer Construction}
\label{ch:lexer-construction}

The rule for an identifier is small; a scanner containing hundreds of overlapping rules is not.
A generator must preserve the meaning of every rule while also deciding which rule wins, where its lexeme ends, and what happens when the next byte belongs to no token.
Chapter~\ref{ch:lexical-analysis} built these decisions directly into control flow.
Here we derive a reusable recognizer from regular expressions, then add the output and rollback policies that turn recognition into lexing.
The construction also supplies the terminal matcher used by the Aurocks generator in Chapter~\ref{ch:aurocks-case-study}.

\section{Separate a language from a token policy}
\label{sec:ch49-language-policy}

Let $\Sigma$ be a finite alphabet and $\Sigma^*$ the finite strings over it.
The empty string is $\epsilon$; it is not the empty language $\emptyset$.
A regular expression denotes a language: union selects either language, concatenation joins one word from each, and star permits zero or more concatenated words.
Character classes are abbreviations for unions of alphabet symbols.
For a byte-oriented scanner, $\Sigma$ can be the $256$ byte values; Unicode scalar values require decoding and a different alphabet interface.
The automata constructions do not decide that interface for us~\cite{hopcroft-2007-automata}.

Consider the expression $a(b\mid c)^*d$.
It accepts \code{ad}, \code{abd}, and \code{abccd}, but not \code{a}, \code{abcdc}, or the empty string.
Acceptance answers a whole-string question.
A scanner instead asks for an accepted prefix at its current position and must distinguish the accepted prefix from the bytes inspected while looking for a longer one.
These are different contracts even when they use the same automaton.

Now combine rules for the keyword \code{if}, identifiers, \code{=}, \code{==}, and decimal integers.
On \code{if2==7}, the identifier wins over the keyword because it consumes more characters.
On \code{if}, both keyword and identifier consume two; a declared rule priority breaks that tie.
Table~\ref{tab:ch49-token-spans} computes the resulting spans.
Priority is consulted only after length, not before it.

\begin{table}[htbp]
\caption{Longest match and then rule priority tokenize six bytes without splitting the identifier.}
\label{tab:ch49-token-spans}
\centering
\begin{tabular}{llll}
\toprule
Input slice & Span & Competing shorter prefix & Chosen kind \\
\midrule
\code{if2} & $[0,3)$ & \code{if} & identifier \\
\code{==} & $[3,5)$ & \code{=} & equality \\
\code{7} & $[5,6)$ & none & integer \\
\bottomrule
\end{tabular}
\end{table}

The output of an accepting state must therefore include a token kind and priority.
Minimizing the combined machine merely as an accept/reject recognizer can merge states that recognize different tokens.
The recognized union language would survive while the compiler's token stream changed.

\section{Build a nondeterministic machine from fragments}
\label{sec:ch49-fragments}

A \term{nondeterministic finite automaton} (NFA) can have several successors for one input symbol and transitions consuming no symbol.
We call the latter epsilon edges.
A convenient compiler represents a fragment by one entry state and a list of still-unconnected outgoing slots.
Patching a slot changes an edge target; it does not duplicate the destination fragment.
Algorithm~\ref{alg:ch49-fragments} describes the recursive construction.
Each recursive call returns fresh states, even when two expression nodes have the same spelling.

\begin{algorithm}[htbp]
\caption{Compile a regular expression into a fragment with patchable exits.}
\label{alg:ch49-fragments}
\begin{pseudocode}
\Require expression tree $r$ over symbols, epsilon, union, concatenation, and star
\Ensure fragment with entry and exit slots, preserving the expression language
\Function{Compile}{$r$}
\If{$r$ is symbol $a$ or epsilon}
  \State create one consuming or epsilon state with an unpatched exit
  \State \Return its entry and its one exit slot
\ElsIf{$r=u\cdot v$}
  \State $U\gets\Call{Compile}{u}$; $V\gets\Call{Compile}{v}$
  \State patch every exit of $U$ to the entry of $V$
  \State \Return entry of $U$, exits of $V$
\ElsIf{$r=u\mid v$}
  \State compile $U,V$; create a split with epsilon edges to both entries
  \State \Return split entry, concatenation of the two exit lists
\Else
  \State $U\gets\Call{Compile}{u}$ where $r=u^*$
  \State create split $s$ with one edge to $U$ and one unpatched edge
  \State patch every exit of $U$ to $s$
  \State \Return $s$, its remaining exit slot
\EndIf
\EndFunction
\State patch the completed fragment's exits to a new accepting state
\end{pseudocode}
\end{algorithm}

Induction on the expression tree explains the construction.
Concatenation requires a successful path through the first fragment before entering the second.
Union introduces a choice without consuming input.
The star split either exits immediately or enters the child and returns to the same choice.
For an expression tree of $R$ nodes, the number of states and edges is $O(R)$.
Construction is also $O(R)$ if exit lists support constant-time concatenation and each slot is patched once; repeatedly copying lists can increase that cost.
Counted repetition must be charged by its expanded size, rather than by the number of digits in its textual bound.

Figure~\ref{fig:ch49-star} shows the star fragment's crucial zero-iteration route.
The accepting state is not placed inside the child: reaching a child exit still permits another iteration.
Nullable children create epsilon cycles, so a matcher needs visited marks even before consuming its first byte.

\begin{figure}[htbp]
\centering
\begin{tikzpicture}[node distance=18mm and 25mm]
\node[basic block] (s) {split};
\node[basic block,right=of s] (u) {child entry};
\node[basic block,right=of u] (v) {child exit};
\node[basic block,below=of u] (out) {next fragment};
\draw[cfg edge] (s)--node[above] {$\epsilon$}(u);
\draw[cfg edge] (u)--node[above,font=\small] {child word}(v);
\draw[back edge] (v) to[bend right=35] node[above] {$\epsilon$}(s);
\draw[cfg edge] (s)|-node[below,font=\small] {zero iterations}(out);
\end{tikzpicture}
\caption{A star fragment loops after a child match and can also leave without consuming input.}
\label{fig:ch49-star}
\end{figure}

\section{Simulate sets, then intern them}
\label{sec:ch49-subsets}

The \term{epsilon closure} $E(X)$ contains every state reachable from $X$ using only epsilon edges, including $X$ itself.
A worklist computes it by marking a state before following its outgoing edges.
For a fixed NFA with $m$ states and $e$ epsilon edges, one closure costs $O(m+e)$ in the worst case with indexed marks.
Clearing the marks also has a cost; generation counters avoid clearing on every small closure but need a wraparound policy.

After consuming symbol $a$, the next active set is
\[
 E\left(\bigcup_{q\in X}\delta(q,a)\right).
\]
One NFA simulation applies this equation at every position.
For bounded-degree expression NFAs, the cost of one length-$L$ match is $O(mL)$ state operations, with $O(m)$ active-set storage.
The same equation constructs a \term{deterministic finite automaton} (DFA): each distinct active set becomes a single named DFA state.
Algorithm~\ref{alg:ch49-subsets} visits only reachable sets.

\begin{algorithm}[htbp]
\caption{Intern reachable NFA subsets as DFA states.}
\label{alg:ch49-subsets}
\begin{pseudocode}
\Require NFA start $s$, alphabet classes $\Sigma$, accepting states with token priorities
\Ensure DFA with the same prefix languages and winning output at every state
\State $X_0\gets E(\{s\})$; intern $X_0$; queue $X_0$
\While{an unprocessed subset $X$ remains}
  \State remove $X$ from the queue
  \State output of $X\gets$ highest-priority accepting token in $X$, or none
  \ForAll{alphabet classes $a\in\Sigma$}
    \State $Y\gets E(\bigcup_{q\in X}\delta(q,a))$
    \If{$Y$ has not been interned}
      \State intern $Y$ and queue it
    \EndIf
    \State set transition $(X,a)$ to the interned identifier of $Y$
  \EndFor
\EndWhile
\end{pseudocode}
\end{algorithm}

The empty subset is the dead state and transitions to itself.
There can be $2^m$ subsets; conversion is not guaranteed to produce a small table.
With $K$ reachable subsets, $A$ alphabet classes, and $O(m)$ transition/closure work per class, a straightforward bounded-degree implementation costs $O(KAm)$ state operations plus interning, and stores $O(Km+KA)$ bits or cells according to representation.
Hashing a bitset is not free; its machine-word length belongs in the interning cost.
Lazy DFA construction computes transitions on demand and can fall back to NFA simulation when a cache budget is exhausted.

For a small independent example, use states $0,1,2,3$ with edges $0\xrightarrow{a}1$, $1\xrightarrow{\epsilon}2$, and $2\xrightarrow{b,c}2$, $2\xrightarrow{d}3$.
Only $3$ accepts.
This compact NFA recognizes $a(b\mid c)^*d$, though it is not the unoptimized fragment output above.
Table~\ref{tab:ch49-subsets} works the conversion completely.

\begin{table}[htbp]
\caption{Five reachable subsets arise before equivalent continuation states are merged.}
\label{tab:ch49-subsets}
\centering
\begin{tabular}{lllllll}
\toprule
Name & Subset & $a$ & $b$ & $c$ & $d$ & Accepts \\
\midrule
$A$ & $\{0\}$ & $B$ & $Z$ & $Z$ & $Z$ & no \\
$B$ & $\{1,2\}$ & $Z$ & $C$ & $C$ & $D$ & no \\
$C$ & $\{2\}$ & $Z$ & $C$ & $C$ & $D$ & no \\
$D$ & $\{3\}$ & $Z$ & $Z$ & $Z$ & $Z$ & yes \\
$Z$ & $\emptyset$ & $Z$ & $Z$ & $Z$ & $Z$ & no \\
\bottomrule
\end{tabular}
\end{table}

\section{Minimize future behavior, including outputs}
\label{sec:ch49-minimize}

Two DFA states are equivalent when every suffix produces the same observable acceptance output from either state.
For a token machine the output includes the winning token kind, not merely an accepting bit.
Partition refinement starts with groups having identical current outputs, then separates states whose transitions reach different groups.
Algorithm~\ref{alg:ch49-refine} uses a simple full-sweep method rather than claiming the bound of a more sophisticated splitter algorithm.

\begin{algorithm}[htbp]
\caption{Refine a complete token DFA until transition signatures agree within blocks.}
\label{alg:ch49-refine}
\begin{pseudocode}
\Require reachable complete DFA, alphabet classes, token output per state
\Ensure quotient DFA preserving every future output
\State partition states by current token output, including none
\Repeat
  \State remember the current partition $P$
  \ForAll{blocks $B$ of $P$}
    \State group states in $B$ by the vector of successor block identifiers
    \State replace $B$ by those nonempty groups
  \EndFor
\Until{no block was split}
\State emit one quotient state per block, with its common output and transitions
\end{pseudocode}
\end{algorithm}

Use a snapshot of the old partition for all signatures in one sweep.
With $K$ states and $A$ classes, there are at most $K-1$ splitting sweeps.
Sorting length-$A$ signatures gives a coarse $O(AK^2\log K)$ comparison-based bound; dictionary-based grouping changes the assumptions and constants.
At stabilization, equivalent blocks stay together under every transition, so induction on suffix length proves preservation.
Conversely, a split records a distinguishing output now or after a finite suffix~\cite{hopcroft-2007-automata}.

In the worked DFA the first output partition is $\{D\}$ and $\{A,B,C,Z\}$.
The symbol $d$ separates $B,C$ from $A,Z$.
Then $a$ separates $A$ from $Z$.
The final blocks are $\{A\},\{B,C\},\{D\},\{Z\}$.
Figure~\ref{fig:ch49-minimal} shows the useful transitions of that four-state result; all omitted transitions go to $Z$.

\begin{figure}[htbp]
\centering
\begin{tikzpicture}[node distance=27mm]
\node[basic block] (a) {$A$};
\node[basic block,right=of a] (b) {$B,C$};
\node[basic block,right=of b,double] (d) {$D$: accept};
\node[basic block,below=of b] (z) {$Z$: dead};
\draw[cfg edge] (a)--node[above] {$a$}(b);
\draw[cfg edge] (b)--node[above] {$d$}(d);
\draw[cfg edge] (b) edge[loop above] node[above] {$b,c$}(b);
\draw[back edge] (d)--node[right,font=\small] {any byte}(z);
\draw[cfg edge] (z) edge[loop below] node[below,font=\small] {any byte}(z);
\end{tikzpicture}
\caption{The subset states $B$ and $C$ have identical futures and become one state; the accepting output remains distinct.}
\label{fig:ch49-minimal}
\end{figure}

\section{Remember acceptance while scanning ahead}
\label{sec:ch49-last-accept}

A DFA can leave an accepting state while investigating a longer candidate.
Failure later does not erase the earlier accepted prefix.
Algorithm~\ref{alg:ch49-munch} keeps a separate saved endpoint and token output.
Token rules must exclude the empty word; otherwise repeated emission can leave the source position unchanged.
End of input is a separate condition, not a byte chosen as a sentinel from the source alphabet.

\begin{algorithm}[htbp]
\caption{Select one longest token while preserving the last accepting endpoint.}
\label{alg:ch49-munch}
\begin{pseudocode}
\Require source $s$ of length $n$, start offset $b<n$, nonnullable token DFA
\Ensure longest token and endpoint, or lexical failure at $b$
\State $q\gets start$; $i\gets b$; $best\gets none$
\While{$i<n$ and the next transition is not dead}
  \State $q\gets\delta(q,s[i])$; $i\gets i+1$
  \If{$q$ has a token output}
    \State $best\gets(i,output(q))$
  \EndIf
\EndWhile
\If{$best=none$}
  \State \Return lexical failure at $b$
\EndIf
\State \Return token spanning $[b,best.end)$, with next cursor $best.end$
\end{pseudocode}
\end{algorithm}

Figure~\ref{fig:ch49-rollback} uses token rules \code{a} and $a^*b$ on \code{aaaa!}.
The scanner remembers acceptance after the first \code{a}, continues looking for \code{b}, then returns to the saved endpoint.
The next token repeats much of that search.
On $a^n$, this implementation can inspect $n+(n-1)+\cdots+1$ bytes, giving quadratic total work despite constant-time DFA transitions.
The linear bound for the direct Mini scanner does not automatically apply to arbitrary maximal-munch specifications.
Bounded rollback, buffering, or memoized failed state/position searches need their own argument~\cite{grune-2008-parsing}.

\begin{figure}[htbp]
\centering
\begin{tikzpicture}[x=17mm,y=10mm]
\foreach \x/\c in {0/a,1/a,2/a,3/a,4/!} {
 \node[draw,minimum width=17mm,minimum height=8mm] at (\x,0) {\ttfamily\c};
}
\draw[cfg edge] (-0.5,1)--node[above,font=\small] {inspect toward a possible \code{b}}(4.5,1);
\draw[back edge] (4.5,-1)--node[below,font=\small] {resume at saved token end}(0.5,-1);
\draw[dashed] (0.5,-1.5)--(0.5,0.6);
\node[font=\small,anchor=west] at (0.7,-2) {accepted span: one \code{a}};
\end{tikzpicture}
\caption{Inspection and consumption have different endpoints; rollback can repeatedly inspect the same suffix.}
\label{fig:ch49-rollback}
\end{figure}

\section{Engineering the character boundary}
\label{sec:ch49-engineering}

A byte span is not automatically a display-column span.
UTF-8 uses variable-length encodings, tabs change display position, and combining characters can share a displayed column.
Keep byte offsets for slicing and separately define decoding and diagnostic coordinates.
An ASCII identifier policy should reject or delegate non-ASCII bytes explicitly; indexing a byte table with a negative signed \code{char} is an implementation bug.

Lexical modes add a finite control component for strings, comments, or interpolation.
Each mode can select a different start state; the token output may switch modes after emission.
Arbitrarily nested block comments require a counter or stack, so that feature is no longer solely a finite-state recognizer.
Do not disguise the extra state as a particularly elaborate regular expression.

The case-study matcher uses NFA simulation and longest prefix matching for each requested terminal separately.
It does not construct one combined token DFA or apply a global keyword priority policy.
That choice keeps the generated runtime small and makes grammar-directed lexical choices possible, but grammar backtracking can invoke the same matcher repeatedly.
Its complexity must include both matching and parser search.

\section{Exercises}
\label{sec:ch49-exercises}
\begin{enumerate}
\item\label{ex:ch49-subsets} Trace \code{acbd} through Table~\ref{tab:ch49-subsets}, including its initial subset. Is \code{acb} accepted?
\item\label{ex:ch49-distinguish} Give a distinguishing suffix for each pair among the four minimized states in Figure~\ref{fig:ch49-minimal}.
\item\label{ex:ch49-cost} For rules \code{a} and $a^*b$, count the successful \code{a}-transition inspections on input $a^5$ across all token attempts.
\item\label{ex:ch49-output} Two accepting states produce different token kinds but have identical outgoing transitions. Explain why accept/reject minimization is insufficient for a lexer.
\end{enumerate}
Checked answers and executable automata checks accompany the front-end examples.

\section{Further reading}
\label{sec:ch49-reading}
Hopcroft, Motwani, and Ullman supply the language-equivalence arguments behind NFA simulation, determinization, and minimization~\cite{hopcroft-2007-automata}.
Grune and Jacobs place lexical recognition alongside parsing strategies, helping distinguish a regular-language matcher from the policy of a complete scanner~\cite{grune-2008-parsing}.
